\chapter*{Résumé}
\addcontentsline{toc}{chapter}{Résumé}

Ce rapport présente les travaux menés durant mon stage de fin d'études au sein
de BCG~X, l'entité de conception et d'ingénierie technologique du Boston
Consulting Group (BCG). Le projet porte sur l'\textit{Energy Transition
Optimizer}, un produit interne d'analyse de pré-faisabilité de projets
\textbf{Power-to-X}, développé puis adapté au cas de
plusieurs clients. Il repose sur un moteur d'optimisation déterministe~: un
projet y est décrit par une topologie (graphe typé d'actifs, de n\oe uds et de
marchés) et des centaines de paramètres, à partir desquels le moteur simule
l'exploitation heure par heure et calcule des indicateurs comme la NPV, l'IRR
et le coût actualisé du produit (LCOX).

Puissant mais réservé aux experts et reposant sur une plateforme fragile, le
produit devait être consolidé et rendu accessible. Mené en équipe, le travail a
modernisé la plateforme et introduit une \textbf{couche conversationnelle
multi-agents gouvernée}~: un agent directeur oriente chaque requête vers des
agents spécialisés (ingestion, conception de topologie, exploration de
scénarios, interprétation des résultats), tandis que toute action modifiant
l'état reste une proposition typée, validée puis soumise à approbation humaine,
le moteur déterministe demeurant la seule autorité numérique.

Trois volets complètent cette couche~:
(i)~l'\textbf{agent d'interprétation des résultats}, un sous-agent en lecture
seule expliquant en langage naturel l'économie des simulations, chaque chiffre
restant tracé jusqu'à un résultat déterministe (aucune valeur inventée)~;
(ii)~la \textbf{couche de visualisation} associée (nuages de points
interprétables et tableau de flux de trésorerie navigable)~; et (iii)~un
\textbf{harnais d'évaluation} « LLM-juge » mesurant l'ancrage factuel et la
solidité financière des explications. Le système allie ainsi la reproductibilité
d'un outil d'analyse d'investissement à l'accessibilité d'une interface assistée
par l'IA.

\noindent\rule[2pt]{\textwidth}{0.5pt}

{\textbf{Mots clés :}}
Power-to-X, pré-faisabilité, optimisation déterministe, systèmes multi-agents,
grands modèles de langage, interprétation de résultats, évaluation LLM-juge,
validation humaine.
\\
\noindent\rule[2pt]{\textwidth}{0.5pt}
